\section*{Chapitre \ref{ch:m2:fx-spot-microstructure} --- Microstructure du change au comptant}

\begin{solution}{exo:m2:fx-spot-microstructure:1}
Le fournisseur a encaissé la demi-fourchette, 100~USD, et perdu le mouvement de deux pips du
milieu en faveur du client, 200~USD~: un \omterm{def:m2:fx-spot-microstructure:markout}{mark-out} de $-\text{USD}~100$, soit $-0.87$
points de base du milieu.
\end{solution}

\begin{solution}{exo:m2:fx-spot-microstructure:2}
$0.01\sqrt{h}$~: 0.05, 0.1 et 0.2 points de base.
\end{solution}

\begin{solution}{exo:m2:fx-spot-microstructure:3}
Si, et comment, les variations du prix dans l'un ou l'autre sens influencent la
décision d'accepter ou de rejeter~; le temps espéré ou typique nécessaire à la décision~; et
la finalité du \omterm{def:m2:fx-spot-microstructure:lastlook}{last look}. Pendant la fenêtre, il ne doit ni utiliser la demande pour recueillir
une information sans intention de négocier, ni négocier sur cette base, en ajustant ses
prix ou en couvrant sa position.
\end{solution}

\begin{solution}{exo:m2:fx-spot-microstructure:4}
$0.05\,\varphi(0) = 0.0199$ points de base à 25 millisecondes, contre 0.0399 à
100~: un quart du \omterm{def:m2:fx-spot-microstructure:lastlook}{temps de rétention}, la moitié du coût.
\end{solution}

\begin{solution}{exo:m2:fx-spot-microstructure:5}
À la baisse~: il abaisse à la fois le cours acheteur et le cours vendeur de 0.1 pip, de sorte que
les clients qui achètent des euros viennent à lui, et ceux qui vendent s'adressent ailleurs. Il paie un dixième de pip sur les transactions
qu'il attire, moins que la fourchette et l'impact de marché lié à la vente de 20~millions d'euros
sur une plateforme, et il ne révèle rien au marché.
\end{solution}

\begin{solution}{exo:m2:fx-spot-microstructure:6}
Il rejette les mouvements au-delà du seuil dans les deux sens, soit une zone deux fois plus large que
celle du contrôle asymétrique. Mais les mouvements rejetés en faveur du client sont compensés par les
mouvements rejetés à son détriment, que le client aurait autrement dû accepter~; le contrôle
asymétrique les conserve et ne refuse que les bons.
\end{solution}

\begin{solution}{exo:m2:fx-spot-microstructure:7}
À 100 millisecondes~: asymétrique, 92.68~\% des requêtes informées et 30.79~\% des requêtes non
informées rejetées, \omterm{def:m2:fx-spot-microstructure:markout}{mark-out} de 0.4491 point de base~; symétrique, 93.46~\% des requêtes informées et
61.76~\% des requêtes non informées rejetées, \omterm{def:m2:fx-spot-microstructure:markout}{mark-out} de 0.3991.
\end{solution}

\begin{solution}{exo:m2:fx-spot-microstructure:8}
Le \omterm{def:m2:fx-spot-microstructure:lastlook}{taux de rejet} ne mesure pas ce qui est rejeté. Un contrôle asymétrique qui rejette
5~\% des requêtes, toutes celles où le cours a bougé en faveur du client, transfère
néanmoins la valeur de ces mouvements~; et un \omterm{def:m2:fx-spot-microstructure:lastlook}{temps de rétention} sans rejet coûte au client
l'option que le fournisseur détient pendant cette durée. Il faut examiner le sens des
mouvements sur les requêtes rejetées et acceptées, les \omterm{def:m2:fx-spot-microstructure:lastlook}{temps de rétention} et les \omterm{def:m2:fx-spot-microstructure:markout}{mark-outs},
et non seulement le nombre de rejets.
\end{solution}

\begin{solution}{pb:m2:fx-spot-microstructure:1}
\textbf{1.} $0.01\sqrt{100} = 0.1$ point de base.
\textbf{2.} $1 - \Phi(0.5) = 30.9\%$.
\textbf{3.} $2(1 - \Phi(0.5)) = 61.7\%$.
\textbf{4.} Asymétrique, $0.1\,\varphi(0.5) = 0.0352$ point de base~; symétrique,
zéro.
\textbf{5.} $0.1/\sqrt{2\pi} = 0.0399$ point de base.
\textbf{6.} $0.0352 \times 10^{-4} \times 2 \times 10^9$, soit environ 7\,041~USD.
\textbf{7.} Environ 1.76~million de dollars.
\textbf{8.} La demi-fourchette coûte 80\,000~USD par jour~; la fenêtre y ajoute 8.8~\%,
de façon cachée.
\textbf{9.} 0.0199 point de base, soit environ 3\,989~USD par jour.
\textbf{10.} Le mouvement pendant le \omterm{def:m2:fx-spot-microstructure:lastlook}{temps de rétention} a un écart-type proportionnel à
$\sqrt h$, et le transfert est proportionnel à cet écart-type.
\textbf{11.} Le sens du mouvement du cours sur les requêtes rejetées (toutes en faveur du
client), les \omterm{def:m2:fx-spot-microstructure:markout}{mark-outs} des requêtes acceptées comparés à ceux des requêtes rejetées, et
des \omterm{def:m2:fx-spot-microstructure:lastlook}{temps de rétention} qui diffèrent selon le sens.
\textbf{12.} Il compte les rejets, sans tenir compte de leur sens~: une faible part de rejets,
tous sur des mouvements favorables, est exactement ce que produit un contrôle asymétrique coûteux.
\textbf{13.} Selon le Principe~17, si et comment les mouvements du prix dans l'un ou l'autre
sens influencent l'acceptation, le \omterm{def:m2:fx-spot-microstructure:lastlook}{temps de rétention} typique et la finalité~; et la confirmation
que la requête n'est pas utilisée pour négocier pendant la fenêtre.
\textbf{14.} Le risque de marché qu'il voulait se décharger, conservé au moment où le
cours évolue contre lui, et le coût de la négociation à nouveau à un moins bon prix.
\textbf{15.} Ses cotations peuvent être plus serrées ou plus larges, et un contrôle symétrique,
court et bien communiqué peut valoir la perte d'exécutions qu'il entraîne.
\textbf{16.} Qu'il a été appliqué de façon large et indiscriminée, sans vérifier
l'existence d'un arbitrage de latence, qu'il n'a rejeté que les transactions non rentables pour
la banque, et que les clients ont reçu des explications trompeuses ou aucune explication.
\textbf{17.} Il ne prend rien en moyenne, mais il refuse tout de même des transactions et
laisse au client le risque~; son équité dépend d'un \omterm{def:m2:fx-spot-microstructure:lastlook}{temps de rétention} court et de la
transparence.
\textbf{18.} Par des écarts plus larges pour les clients aux \omterm{def:m2:fx-spot-microstructure:markout}{mark-outs} défavorables, des tailles plus petites,
des paliers, des mises à jour de prix plus rapides, et en refusant de diffuser des cotations aux flux toxiques.
\textbf{19.} Résultat nommé~: \emph{fenêtre asymétrique} de 100 millisecondes, avec
un seuil de 0.05 point de base, coûte 0.0352 point de base par requête, soit environ
1.76~million de dollars par an pour 2~milliards de dollars de requêtes par jour.
\textbf{20.} La valeur des mouvements de prix en sa faveur que le fournisseur refuse
d'honorer, à laquelle s'ajoute le risque de chaque transaction rejetée.
\end{solution}

\begin{solution}{iq:m2:fx-spot-microstructure:1}
Une courte fenêtre qui suit l'arrivée d'une requête de transaction, pendant laquelle le
fournisseur peut l'accepter ou la rejeter au prix coté. Les fournisseurs l'utilisent pour protéger
les prix diffusés contre l'arbitrage sur prix périmés et sur latence, pour vérifier le crédit et la
validité, et ainsi coter plus serré pour tous~; le coût retombe sur les clients dont les
requêtes sont rejetées.
\iqlookfor{la définition, l'argument de protection et son coût.}
\end{solution}

\begin{solution}{iq:m2:fx-spot-microstructure:2}
Enregistrer pour chaque requête le milieu au moment de la requête, au moment de la décision et après celle-ci.
Comparer la distribution du mouvement pendant le \omterm{def:m2:fx-spot-microstructure:lastlook}{temps de rétention} pour les requêtes rejetées et
acceptées~: avec un contrôle symétrique, les rejets surviennent sur de grands mouvements dans les deux
sens~; avec un contrôle asymétrique, seulement sur des mouvements en faveur du client. Tester si la
probabilité de rejet dépend du signe du mouvement, et si les \omterm{def:m2:fx-spot-microstructure:lastlook}{temps de rétention}
diffèrent selon le sens.
\iqlookfor{conditionnement selon le sens du mouvement.}
\end{solution}

\begin{solution}{iq:m2:fx-spot-microstructure:3}
Le profit du fournisseur sur une transaction marquée au milieu à un horizon ultérieur~: la
demi-fourchette moins le mouvement favorable au client. Moyenné par client sur des horizons
allant de la milliseconde à la minute, il révèle qui trade avant les mouvements. Le classer par paliers~:
cours plus serrés pour les clients aux \omterm{def:m2:fx-spot-microstructure:markout}{mark-outs} positifs, cours plus larges ou plus lents pour
les autres, réestimés régulièrement avec des intervalles de confiance.
\iqlookfor{la définition et une procédure de classement raisonnable.}
\end{solution}

\begin{solution}{iq:m2:fx-spot-microstructure:4}
Il dépend du flux attendu et des limites de risque. Si des clients sont susceptibles
d'acheter des euros prochainement, décaler le prix pour les attirer et internaliser~; si la position est
importante par rapport à la limite ou si le marché est rapide, couvrir une partie sur une plateforme
ou avec un algorithme d'exécution. Attendre sans décalage revient à prendre position sur le prix,
ce que doit faire un teneur de marché seulement de manière délibérée.
\iqlookfor{inventaire, flux attendu et limites, pas une règle unique.}
\end{solution}

\begin{solution}{iq:m2:fx-spot-microstructure:5}
Transparence et information sur le fonctionnement du \omterm{def:m2:fx-spot-microstructure:lastlook}{last look}, y compris sa symétrie et son
calendrier~; ne l'utiliser que comme contrôle de validité et de prix~; aucune collecte d'informations
et aucune transaction sur la requête pendant la fenêtre. Il n'interdit pas les contrôles asymétriques
et ne fixe pas les \omterm{def:m2:fx-spot-microstructure:lastlook}{temps de rétention}~; il n'a pas force de loi~: l'adhésion est volontaire,
par des engagements déclaratifs.
\iqlookfor{le fond du principe 17 et ses limites.}
\end{solution}

\begin{solution}{iq:m2:fx-spot-microstructure:6}
Un cœur de calcul unique calcule le milieu à partir des flux des plateformes~; les écarts et les décalages
par palier sont appliqués dans une table précalculée~; les cotations sont publiées par une couche de diffusion
avec des limiteurs de débit par client. Les requêtes arrivent sur un chemin distinct~; le contrôle compare
le milieu courant au milieu à l'instant de la cotation, stocké avec chaque identifiant de cotation,
à un seuil symétrique, après un \omterm{def:m2:fx-spot-microstructure:lastlook}{temps de rétention} fixe et court~; la décision et ses données d'entrée
sont journalisées pour une analyse ultérieure. Aucune allocation sur le chemin critique, des files sans verrou
et un test de rejeu sur des données de marché enregistrées.
\iqlookfor{séparation des chemins de tarification et d'acceptation, logique symétrique, journalisation
pour l'audit.}
\end{solution}
